\section{Teleporters Path}
\label{sec:cses-teleporters-path}

\problemstatement{Given a \emph{directed} graph of $n$ nodes and $m$ edges,
find a path from node $1$ to node $n$ that uses \emph{every edge exactly
once} (a directed Eulerian path), or output \texttt{IMPOSSIBLE}.}

\begin{patternbox}
A directed \textbf{Eulerian path} from $1$ to $n$ exists iff node $1$ has
one more outgoing than incoming edge, node $n$ one more incoming than
outgoing, every other node has equal in- and out-degree, and all edges are
connected. \textbf{Hierholzer's algorithm} then constructs it.
\end{patternbox}

\subsection*{Degree imbalance pins the endpoints}

Check the directed degree conditions: $\textit{out}[1]-\textit{in}[1]=1$,
$\textit{in}[n]-\textit{out}[n]=1$, and $\textit{in}=\textit{out}$ for all
other nodes; also confirm every edge is reachable so the trail is a single
path. Then run Hierholzer from node $1$: follow unused out-edges via a
per-node pointer, pushing onto a stack; when a node has no unused out-edges,
pop it to the path. Reversing the popped nodes gives the Eulerian path,
which necessarily ends at $n$.

\begin{ideabox}[The +1/-1 Imbalance Fixes Start and End]
Equal in/out degree everywhere would force a circuit; a single node with an
extra out-edge must be the start and a single node with an extra in-edge the
end. Hierholzer respects this automatically, terminating the trail at $n$.
\end{ideabox}

\subsection*{Algorithm}

\begin{enumerate}
  \item Compute in/out degrees; verify the $1\to n$ Eulerian-path
        conditions, else \texttt{IMPOSSIBLE}.
  \item Hierholzer from node $1$ over directed out-edges with a per-node
        pointer.
  \item Reverse the popped nodes; if all $m$ edges are used and it ends at
        $n$, print it.
\end{enumerate}

\subsection*{Dry run}

\dryrun{$1\!\to\!2$, $2\!\to\!3$, $1\!\to\!3$? (degrees must fit).}

\begin{itemize}
  \item With $1\!\to\!2$, $2\!\to\!3$: $\textit{out}[1]{-}\textit{in}[1]{=}
        1$, $\textit{in}[3]{-}\textit{out}[3]{=}1$, node $2$ balanced.
  \item Path $1\to2\to3$ uses both edges.
\end{itemize}

\subsection*{C++ implementation}

\begin{lstlisting}
#include <bits/stdc++.h>
using namespace std;

int main() {
    int n, m;
    cin >> n >> m;
    vector<vector<int>> adj(n + 1);
    vector<int> indeg(n + 1, 0), outdeg(n + 1, 0);
    for (int i = 0; i < m; i++) {
        int a, b; cin >> a >> b;
        adj[a].push_back(b);
        outdeg[a]++; indeg[b]++;
    }

    // Eulerian path 1 -> n conditions
    bool ok = (outdeg[1] - indeg[1] == 1) && (indeg[n] - outdeg[n] == 1);
    for (int i = 1; i <= n && ok; i++)
        if (i != 1 && i != n && indeg[i] != outdeg[i]) ok = false;
    if (!ok) { cout << "IMPOSSIBLE\n"; return 0; }

    vector<int> ptr(n + 1, 0), path;
    stack<int> st;
    st.push(1);
    while (!st.empty()) {
        int u = st.top();
        if (ptr[u] < (int)adj[u].size()) {
            int v = adj[u][ptr[u]++];
            st.push(v);
        } else { path.push_back(u); st.pop(); }
    }
    reverse(path.begin(), path.end());

    if ((int)path.size() != m + 1 || path.front() != 1 || path.back() != n) {
        cout << "IMPOSSIBLE\n"; return 0;
    }
    for (int v : path) cout << v << " ";
    cout << "\n";
    return 0;
}
\end{lstlisting}

\begin{complexitybox}
\textbf{Time Complexity.} $O(n+m)$. \textbf{Space Complexity.} $O(n+m)$.
\end{complexitybox}

\begin{takeawaybox}
Directed Eulerian paths need the in/out degree imbalance to name the
endpoints, then Hierholzer. The final check -- all edges used and the trail
runs $1\to n$ -- guards against disconnected edge sets that pass the local
degree test.
\end{takeawaybox}
